\chapter[Islam: Faith, Community, Knowledge]{Islam: Revelation, Community, and Knowledge}
\section{A Gold Coin Without a King's Face}
First, pick up a gold coin.

In the coin displays of a major museum in London, Berlin, or elsewhere, you may find it under ``Early Islam'': about two centimetres across, thin, heavy, high-purity gold, minted in late-seventh-century Damascus. It differs from its neighbours. Roman gold shows an emperor's profile; Persian silver, a king and a Zoroastrian altar. This coin has no person, animal, or image, only carefully arranged rings of Arabic writing. Its central message is one sentence: there is no god but God.

Money is society's most widely circulated newspaper. In antiquity's low-literacy world, a coin was an announcement everyone handled. An emperor's face declared what power looked like and demanded recognition. The late-seventh-century reform led by Umayyad caliph Abd al-Malik was more radical. He abolished the Byzantine- and Persian-image coins Arabs had previously used and replaced them with new money bearing no face, only a statement.

What statement? A profession of faith, which every merchant, farmer, tax collector, and soldier handling the coin had to transmit by hand.

Strange, is it not? Minting a declaration makes every exchange a broadcast. What community considers a sentence more important than a leader's face?

This coin is our key. Behind it stands a civilisation beginning with monotheistic revelation, uniting Arabia within decades, extending from Spain to Central Asia within a century, and becoming a centre of knowledge and commerce within several centuries. It is Islam. We will examine not a checklist of religious facts but three intertwined things: revelation, what believers think God said; community, how people organised around it; and a world of knowledge, the minds that community nurtured.

Leave the coin on the table. Go back a century before it appeared, to an Arabian Peninsula where none of this had happened.

\section{A Night in Mecca, 610 CE}
First, a map.

Two giants occupied western Asia in the early seventh century: Byzantium, the Eastern Roman Empire, to the west, and Sasanian Persia to the east. More than twenty years of major warfare had exhausted their treasuries and populations. In the space beneath them lay Arabia, vast deserts and oases neither much valued, treated as a buffer zone by Romans and Persians.

Arabs inhabited the peninsula. Oasis towns along the Red Sea conducted transit trade, while Bedouin pastoralists followed water and grazing through the desert. To understand this world, understand tribe. There was no state, police, or court. Kinship was your entire personal insurance: if someone attacked you, your tribe owed you vengeance. Bloody as this sounds, it was an unwritten law of deterrence. The certainty of retaliation made attackers think twice.

Mecca was an exception rising from this desert: an oasis city with the famous Zamzam spring and a cubic stone sanctuary, the Kaaba. It housed idols belonging to the peninsula's tribes, more than three hundred according to tradition. People came throughout the year to visit their gods and trade. Mecca's Quraysh tribe supplied water, maintained order, and took a share of commerce. The tribes also recognised sacred pilgrimage months in which fighting was forbidden. Mecca became a desert free-trade and ceasefire zone. Guarding stone gods, the Quraysh effectively managed the peninsula's peace and logistics.

Muhammad was born here around 570 CE. His life began badly: his father died before his birth, his mother when he was six, and his uncle raised him without wealth or power. As a young man he travelled with caravans and served as a trade agent. Tradition says his reputation earned him the name al-Amin, the trustworthy. Around twenty-five, he married his employer Khadija, a wealthy, independent widow around forty.

In 610 he was forty. Islamic tradition records that during Ramadan he went alone, as usual, to reflect in a cave on Mount Hira outside the city. One night he underwent a terrifying experience: a mighty presence commanded, ``Recite!'' He replied that he could not read. The command came again: ``Recite!''

Trembling, he ran home and asked his wife to wrap him in a cloak. Hearing his account, Khadija spoke words generations of Muslims remembered: God would not disgrace someone who supported relatives, cared for the poor, and helped those in distress.

Muhammad believed the voice belonged to an angel sent by God. Revelation continued for twenty-three years until his death, recorded and memorised to form the Quran. He understood himself as the last prophet in a chain including Abraham and Moses, familiar to Jews, and Jesus, familiar to Christians. He was its seal, its final link.

Mecca's night was quiet. Nobody knew the peninsula's history had just been divided in two.

\section{One Person, One Message, One Peninsula}
Do not take sides yet. First see the conflict.

The message Muhammad brought from the cave sounded simple, but every part undermined Mecca's arrangements.

First, one God. Not three hundred gods, but one; the Kaaba's stone images were false. That amounted to saying Mecca's economic lifeline rested on deception.

Second, final judgement. After death, people would account for how they acquired wealth and treated orphans and the poor. The rich were being told their ledgers would face a court.

Third, kinship was not the highest tribunal. Before God, tribal honour counted for nothing; slave and chief faced the same scales.

The fourth point was most consequential. It would acquire a powerful name: umma, community, bound not by blood but by shared commitment. Belonging depended not on one's father but on what one acknowledged.

Now you can see why Meccan elites lost sleep. This was not a theological disagreement about one more or fewer gods, but a demolition notice for an entire order.

Our real question emerges, neither romantic nor mysterious:

\textbf{How could a minority without authority, an army, or hereditary standing remake the peninsula within a generation, then grow into a world civilisation encompassing Persians, Turks, Indians, Malays, and West Africans?}

By swords, the message itself, or organisational genius? Establish one rule before continuing: distinguish \textbf{what believers believe}, that Muslims understand the Quran as divine revelation; \textbf{what tradition claims}, the accounts of scholars and histories across ages; and \textbf{what historical evidence shows}, what documents, coins, archaeology, and outside testimony establish. They often overlap but are not identical. Understanding a faith does not require believing it; explaining its rise does not defend everything done in its name.

\section{Whose Mecca? Four Voices}
Between 610 and 622, Mecca argued over Muhammad's message for twelve years. Hear the sides.

\textbf{First voice: Khadija.} His first believer was not a man but his wife, around fifteen years older, his employer, financier, and spiritual support. Her wealth sustained the small group through the hardest years of preaching. Notice: the movement's first adherent was a businesswoman controlling her own property.

\textbf{Second voice: Bilal.} An Abyssinian slave, from around present-day Ethiopia, he was among the earliest converts. Islamic tradition says his master dragged him into the midday desert, pinned a great stone on his chest, and demanded renunciation. He repeated only, ``One, One.'' Early believers later raised money for his freedom. In Medina, the former slave became Islam's first muezzin, calling the city to prayer five times daily. A slave's voice became the community's alarm clock. You can see why the poor, enslaved, and unprotected first welcomed a message saying their scales were the same as everyone else's.

\textbf{Third voice: Mecca's aristocrats.} Often reduced to stock villains, they deserve their full case. First, Kaaba pilgrimage sustained Mecca's economy and the peninsula's truce. Hundreds of tribes traded peacefully through a shared arrangement of many gods and sacred months. Smashing idols meant smashing livelihoods and order. Second, ancestral ways should not lightly be abandoned, a basic ancient ethic, not simply ignorance. Third, they had legitimate doubts: why should an orphan without army or royal blood speak for God? Fourth, tradition says they initially offered bargains rather than violence: wealth, rank, even alternating worship as compromise. Muhammad refused. These reasons show a functioning economic, security, and moral order being defended, not mere stupidity. Yet complete reasons do not guarantee correct conclusions. They later used beatings, boycott, and threats of assassination; slaves and the poor suffered most. Those facts remain.

\textbf{Fourth voice: Medina.} More than four hundred kilometres north lay Yathrib, an oasis whose two Arab tribes were exhausted by years of fighting. Several Jewish tribes also lived there, farming and practising crafts, with scriptures of their own. In 622, people invited Muhammad as a mutually trusted, neutral mediator. He and his followers left Mecca by night in the Hijra, the migration. That year later became year one of the Islamic calendar. Notice the choice: not the prophet's birth or first revelation but \textbf{the community's establishment}. Yathrib became Medina, the Prophet's City.

There Muhammad became lawgiver and commander as well as preacher. A document known as the Constitution of Medina describes the arrangement: Meccan migrants, local tribes, and Jewish tribes formed one umma, defended together, retained their religions and internal affairs, and referred disputes to God and Muhammad. Most historians regard it as genuinely reflecting contemporary arrangements. Watch the wording: umma's birth certificate described a multireligious urban alliance.

What followed was not gentle. Medina and Mecca fought for years; alliances between the new authority and Jewish tribes broke down in succession. Some were expelled. In the harshest case, the men of Banu Qurayza were executed after conflict. This is not sensationalism but a reminder that from the beginning the community faced an unresolved question: where does the boundary of us lie, and what happens to those outside it?

In 630, Muhammad returned to Mecca with an army and almost no bloodshed. He removed the idols, preserved the Kaaba, and pardoned former persecutors, including the prominent Quraysh leader Abu Sufyan. He died in 632.

That afternoon the community faced every departing founder's question: who succeeds? Different answers divided Islam into major traditions that endure today.

\section{One Succession Problem, Three Answers}
Muhammad died without a surviving son or a written succession arrangement, or at least the parties remembered his arrangements differently.

That day, Medina's local Helpers gathered beneath a shelter to propose one of their own. Meccan emigrants joined them. After argument, those present pledged allegiance to Abu Bakr, Muhammad's close friend and an early convert. He became the first caliph, literally successor.

Others were silent but later became a great historical voice: Muhammad's cousin and son-in-law Ali and members of the prophetic family. They thought leadership belonged in that family; Ali was kin and an early believer. Defeated then but not extinguished, they became the shiat Ali, party of Ali, or Shia.

The following thirty years seem fast-forwarded. Under the second caliph, Umar, Arab armies took Syria, Egypt, and all Sasanian Persia within little more than a decade. The giant that had stood beside Byzantium for centuries was swallowed. In less than a generation, Medina's small city-state became an empire across three continents.

Then succession returned to collect its debt. The third caliph, Uthman, was assassinated in 656. Ali became the fourth, but Syria's governor Muawiya, Uthman's relative, refused recognition. Warfare reached stalemate and arbitration. This enraged radicals in Ali's camp: how could God's cause become human bargaining? They withdrew, becoming the Kharijites. Their third answer chose leadership by piety alone, not descent or negotiation. A Kharijite killed Ali in 661. Muawiya established the Umayyad dynasty in Damascus, making the caliphate hereditary.

In 680, Muawiya's son Yazid succeeded. Ali's son Husayn, Muhammad's grandson, refused allegiance and travelled north with dozens of relatives and followers. Umayyad troops surrounded them at Karbala in Iraq and killed them all. It was the tenth day of the first Islamic month, Ashura. Millions of Shia Muslims still mourn Husayn annually from Tehran to Beirut. A catastrophe over thirteen centuries old remains a living wound.

Three answers had formed. The majority held that respected community members should choose leaders through consultation, guided by the Prophet's practice, the sunna: Sunnis, now eighty to ninety percent of Muslims. Shia hold that leadership belongs to the prophetic family and successive imams possess special spiritual authority. Today they form somewhat over ten percent, concentrated in Iran, Iraq, Azerbaijan, and Bahrain. Kharijites long ago became marginal as a sect, but the impulse to demand piety above everything repeatedly returns in history.

Remember two honest qualifications. First, the split began politically, over power, not theology; theological differences developed over centuries. Second, neither major tradition is a single block. Sunnis have many schools, and Shia several branches. Reducing a tradition to one voice is a common misreading.

\section[Thinking Bridge: Three Layers of Reading]{Thinking Bridge: Scripture, Tradition, and History}
You now hold a cave, the command to recite, a slave, a constitution, an argument under a shelter, and Karbala. Are they reliable? Who tells them? Here is a portable tool for any religious tradition: read in three layers.

\textbf{First, what scripture says: the textual layer.} The Quran's 114 chapters are available for anyone to check, with tens of thousands in each generation able to recite the whole. This layer is least disputed, not because meaning is undisputed but because the written wording can be examined.

\textbf{Second, how tradition remembers: the traditional layer.} The cave's details, Bilal's stone, and Khadija's words come from hadith and prophetic biographies, transmitted orally before being written. Muslim scholars scrutinised them intensely. They developed methods resembling detailed population investigations, checking chains in which A heard B, B heard C, whether each transmitter was honest and could have met the preceding person, then grading reports. Scepticism towards one's own tradition became a formal discipline. It also acknowledged gaps: falsely attributed prophetic reports exist.

\textbf{Third, what historical evidence shows: the historical layer.} This uses outside corroboration rather than believers' memory alone: non-Arab accounts, coins, inscriptions, and archaeology. It confirms that a seventh-century movement centred on Arabic scripture rose in Arabia and built a great empire within decades. Our opening coin is a material witness at this level.

Practise with an easily misjudged example. A supposed prophetic saying widely quoted in Chinese urges seeking knowledge even as far as China. You may have met it in books or speeches. Three-layer reading finds it absent from the Quran, doubtful in hadith scholarship because its transmission chain is weak, and recorded much later than the Prophet. It is probably a well-intentioned later attribution rather than Muhammad's words. The example works both ways: do not trust a saying merely because it circulates widely, and notice that this tradition's own textual scholars, not outsiders, first exposed its weakness.

The method applies beyond religion. Family recollections, viral screenshots, and textbook stories invite the same questions: what did the original text say, what did collective memory add, and what can independent evidence confirm?

\section[Recite: Scripture in a Few Lines]{Recite: A Scripture's Character in a Few Lines}
Read a short primary passage. Muslims believe the opening verses of Quran chapter 96 began revelation. Rendered here from Ma Jian's Chinese translation:

\begin{quote}
Recite in the name of your Creator, who created the human being from a blood clot. Recite: your Lord is most generous, who taught by the pen, who taught humanity what it did not know.
\end{quote}
Read slowly. The passage contains the scripture's character.

Its first command is to recite; Quran itself means recitation. For Muslims, the Quran is first a sound, not a book: memorised in prayer, chanted on broadcasts, learned by countless children before literacy. A religion beginning with a command to read later built world-class libraries. Remember that connection.

Notice the manner: in the Creator's name. Reading is not a private hobby; one first declares knowledge's source. The pen is explicitly sacred, the instrument by which God teaches. Seventh-century Arabia prized oral poetry, and writing belonged to a minority. Revelation elevated writing to a central divine gift. A culture of ears would grow a culture of hands.

The third passage worth reading closely is chapter 49, verse 13, also from Ma Jian's version:

\begin{quote}
O people, I created you from a man and a woman, and made you peoples and tribes so that you might know one another. In God's sight, the most honoured among you is the most reverent.
\end{quote}
In tribal Arabia this sounded like thunder. Peoples and tribes existed not to conquer one another but to become acquainted. Rank depended not on descent, tribal society's foundation, but on reverence, an inward quality available to any origin. This was what people like Bilal heard.

One more line will recur: the beginning of chapter 2, verse 256, ``There is no compulsion in religion.'' Next, compare it with history.

\section[Fast Conquest, Slow Conversion]{Comparing Explanations: Fast Conquest, Slow Conversion}
The facts: fewer than a hundred years separate Muhammad's death in 632 from Umayyad expansion to Spain and Central Asia in the early eighth century, an exceptional historical speed. Equally important, conversion in conquered Syria, Egypt, Persia, and North Africa was much slower, taking two or three centuries or more and never becoming complete. Ancient Coptic Christian communities remain in Egypt, and many Christians still inhabit Lebanon's mountains.

One set of facts, three explanations.

\textbf{Explanation one: military strength and timing.} Arab armies were mobile, brave, and motivated; Byzantium and Persia were exhausted by decades of warfare, empty treasuries, and plague. Sectarian tensions between Syria and Egypt and imperial centres reduced willingness to resist a new ruler to the death. Conquest was indeed largely military. Its limit is obvious: history has many rapid conquests that disappeared as quickly as steppe cavalry. Why did these lands not simply revert, instead retaining the conquerors' language and religion? Military power does not explain lasting attachment.

\textbf{Explanation two: organisation and commitment.} The umma offered an unusual identity based on shared commitment and law rather than tribe or ethnicity. The new regime gave conquered ``People of the Book'', including Jews and Christians, protected status: a special tax, jizya, in exchange for security of person, property, and religious autonomy. Neither romanticise nor demonise this immediately. By modern standards it was unequal toleration, legally second-class status. Yet a seventh-century conqueror permitting organised non-Muslim communities without compulsory conversion was less common than convert-or-die conquest. The verse against compulsion entered institutional history here. Not always observed, it nonetheless became an arguable, invocable standard. This explains durability but risks self-congratulation that understates genuine discrimination.

\textbf{Explanation three: economy and synthesis.} The conquerors were few and initially lived in garrison cities, governing through existing bureaucrats, tax officials, and scholars. Persians continued Persian paperwork; Syrian Christians kept accounts. The superstructure was Arab, its foundations Persian, Greek, and Syrian. Arabic took one or two centuries to become administrative language, while Persian later revived in the east as a new literary language. Rather than Arabs remaking old civilisations, perhaps those civilisations absorbed their conquerors and jointly produced something new. This best explains later cultural flourishing but softens conquest's violence: cities were stormed, taxes real, and captives' fates real.

Each explanation holds a piece. Together: military conquest was retained by an open community design, then nourished by ancient civilisations' accumulated resources. See what this mixed inheritance produced.

\section{Baghdad's Translation Production Line}
In 762, an Abbasid caliph drew a perfect circle beside the Tigris and founded Baghdad. For the following two or three centuries, it was among Earth's densest concentrations of knowledge.

Practical interests began the process: caliphs needed medicine, astrology, agriculture, and statecraft. In the ninth century, al-Mamun established the House of Wisdom and funded systematic translation of Greek, Persian, and Indian learning into Arabic. Look at the workers. Leading translator Hunayn ibn Ishaq was a Nestorian Christian physician. With his son and pupils, he translated many medical works of Galen and Hippocrates, often first into Syriac, then Arabic. Persian scholars brought Sasanian court learning and Indian numerals; Harran's Sabians contributed mathematics and astronomy. Later accounts say translators earned their manuscripts' weight in gold. The amount may be exaggerated, but the direction is right: knowledge was hard currency.

Now the output. Several names remain close to your daily life:

\begin{itemize}
\item Al-Khwarizmi, a ninth-century scholar of Persian origin. His short book on solving equations included al-jabr, restoration, in its title, giving algebra its name. His Latinised name, Algoritmi, gave us algorithm.
\item Ibn Sina, known in the West as Avicenna, a Central Asian around the eleventh century. His Canon of Medicine, translated into Latin, served European medical schools for centuries, into the seventeenth century.
\item Ibn al-Haytham, an eleventh-century optical scholar whose major work was completed in Cairo. His method of questioning ancient authorities, experimenting, and proving geometrically is regarded by many historians as a precursor of modern scientific method.
\item Al-Biruni, a later compatriot of al-Khwarizmi, estimated Earth's radius using a mountain and geometry, and wrote a major, careful study of Indian religion and science.
\item Ibn Rushd, Averroes in the West, a twelfth-century Spanish judge and philosopher who commented on Aristotle line by line. His Latin versions stirred European universities, making Averroists a prominent scholastic camp.
\item Do not forget Maimonides, a Jewish philosopher contemporary with Ibn Rushd and also born in Cordoba. He wrote one of Judaism's greatest works in Arabic.
\end{itemize}
Count the ingredients: Persians, Central Asians, Spaniards, North Africans, Muslims, Christians, and Jews. Fix this point firmly: \textbf{the Islamic world has never belonged exclusively to Arabs or to Muslims}. It was a republic of learning whose shared language was Arabic, later Persian too, and contributors of every background counted as authors. Essential hardware came from the east: papermaking arrived through Central Asian Samarkand in the mid-eighth century, replacing costly parchment with inexpensive paper. Without paper, no such translation production line.

Knowledge followed trade routes. Arab dinars and Persian dirhams reached the Baltic; thousands of travelled silver dirhams have been excavated on Sweden's Gotland. Westwards, ships reached Guangzhou and Quanzhou, whose Tang and Song ports housed Arab and Persian merchant quarters. Quanzhou's Qingjing Mosque still stands among China's oldest surviving mosques. Southwards, caravans crossed the Sahara, exchanging salt for gold, and Timbuktu became a city of books. In the fourteenth century, Mali's ruler Mansa Musa made the Meccan pilgrimage, distributing so much gold that contemporary accounts say Cairo's gold price took years to recover.

\section[Further Challenge: Ritual Makes Community]{Further Challenge: Durkheim's Glue and Ritual Community}
Bring in the theoretical tool. Until now we asked what the Islamic world is. Now ask more deeply: \textbf{what keeps a community from falling apart?}

In the early twentieth century, French sociologist Émile Durkheim proposed in The Elementary Forms of Religious Life that religion's central function is not belief in gods but society periodically remaking itself. People gather, perform shared actions, and sing together, bodily renewing the feeling of us. For him, sacred feeling is society's own force experienced as divine. Ritual is not belief's decoration but community's power supply.

Through these spectacles, Islam's Five Pillars look like a precisely designed community engine:

\begin{itemize}
\item \textbf{Profession of faith}: there is no god but God, and Muhammad is God's messenger. Conversion requires no baptism or examination, only sincere public recitation: a low threshold and strong consistency.
\item \textbf{Prayer}: five daily prayers, all facing Mecca's Kaaba. Imagine the arrangement: from Jakarta to Casablanca, more than a billion people point their bodies towards one place five times daily, connecting personal clocks to the community's clock.
\item \textbf{Fasting}: during Ramadan, no food or drink from sunrise to sunset. Synchronised discipline becomes a festival: people hunger together and break the fast together, rich and poor sharing thirst.
\item \textbf{Almsgiving}: zakat assigns a fixed share of savings, commonly one fortieth, to the poor. Notice the design: redistribution is not left to spontaneous moral goodwill but incorporated into worship. Giving is a duty, not merely charity.
\item \textbf{Pilgrimage}: those able to travel safely should visit Mecca at least once. Around two million attend annually today, all wearing identical white pilgrimage clothing, making king and taxi driver indistinguishable as they circle the same sanctuary. Durkheim would point and say society is recharging itself.
\end{itemize}
The explanatory power is striking: without shared blood, territory, or an incumbent ruler, rhythms repeatedly bind strangers into an umma.

But the theory has at least two gaps.

First, it cannot explain inward experience. Durkheim's ritual concerns us, but Sufis have continually sought the direct relation between I and God. Their name is said to derive from suf, the coarse wool they wore. Chanting, meditation, poetry, and dance pursue union with the sacred. The thirteenth-century Rumi, born in present-day Afghanistan and celebrated in Konya, Turkey, composed the enduring Persian Masnavi; the order he founded is famous for whirling. Sufi teachers also led Islam's spread: Southeast Asian islands, Indian villages, and West African towns often entered the umma through teachers and merchants rather than armies. Sufism reveals internal tensions too. The tenth-century al-Hallaj was executed in Baghdad for declaring ``I am the Truth''. Collective ritual and individual experience pull against one another throughout Islamic history.

Second, it cannot explain doctrinal conflict. People sharing rituals still draw swords over theology. In a famous ninth-century dispute, Mutazilites defended reason as faith's interpretive tool and held the Quran created; traditionalists insisted it was God's eternal speech. Caliph al-Mamun supported the rationalists and instituted the Mihna, a doctrinal inquisition of scholars. Jurist Ibn Hanbal endured torture rather than concede, and this coercive state theology ultimately failed. Read the irony carefully: a community whose scripture rejects religious compulsion could have a government imposing theological unity, resisted by a scholar willing to suffer. Afterwards a broad division developed: caliphs governed swords, ulama interpretation. Church and state intertwined without merging.

Speaking of ulama, dismantle a current misunderstanding. News often depicts sharia as a code of cruel punishments. Historically, its root meaning is a path to water. Rather than one written code, it is a living argumentative tradition: jurists reason from Quran and hadith through consensus and analogy about marriage, inheritance, contracts, commerce, and religious duties. Four major Sunni legal schools coexist, alongside Shia schools, and can reach different conclusions in the same case. It resembles a growing forest of precedents more than an immovable stone. Historically, family affairs and business occupied it more than criminal punishment.

\section{Echoes Today}
This fourteen-hundred-year-old story continues quietly in your city.

First, dismantle three linked mistakes. Muslims are not synonymous with Arabs. Nearly two billion Muslims inhabit the globe; the largest national population is in Indonesia, over four thousand kilometres from Mecca, where merchants and Sufi teachers brought Islam peacefully over centuries. Javanese shadow theatre still stages Indian epics while performers begin with the profession of faith. Nor do all Muslims speak Arabic. Iranians use Persian, Turks Turkish, Pakistanis Urdu: Arabic in worship, diverse languages in daily life. Arabs are a minority of Muslims. Finally, Arabs are not all Muslim. Egypt's Copts and Lebanese and Syrian Christian communities predate Islam and likewise speak Arabic as their mother tongue.

Look at China. Islam arrived over thirteen hundred years ago. Xi'an's Great Mosque in Huajue Lane has upturned roofs and bracketed eaves like a Chinese courtyard; only deeper inside do you notice its west-facing prayer hall. Hui people speak Chinese and explain their faith through Chinese ethical vocabulary. Zheng He, voyager to the western seas, came from a Yunnan Muslim family. This civilisation is not merely across the map: it is your classmate, the noodle shop downstairs, and the Arabic inscription in your provincial museum.

Then the unavoidable contemporary controversy. Separate four matters. Verified facts: organisations invoking Islam have committed massacres in recent decades; Muslims themselves are extremism's largest victim group, and mainstream Islamic scholarly institutions have publicly condemned those organisations. Competing explanations: one attributes violence to doctrine, another to politics, including colonial legacies, war, occupation, and failed government using religion for mobilisation. The author's judgement favours the second because it covers more facts: the same texts do not produce violence in most places and periods, while violence follows warfare and collapsing regimes. The open question is how responsible a tradition is for crimes committed in its name. That question applies to every tradition, including yours.

Finally, a quieter echo. Ninth-century Baghdad translated Greek, Persian, and Sanskrit into Arabic because knowledge's capital stood there. Today papers are written in English because its capital stands here. No language is inherently noble; one temporarily inhabits the centre of knowledge, which has moved and will move again. As you study English, consider whether you are learning a language or acquiring entry to the world of knowledge behind it.

\section{For You: One Statement, One Community}
Return to the coin.

A rising community reminted money, replacing a king's face with a sentence. The face had been an obvious adhesive: recognise it and know who commands. The sentence replaced it with commitment: do you acknowledge this declaration?

The umma's ambition hides in that gesture: a human community bound not by blood, appearance, or birthplace but by shared principles. It truly placed Ethiopian slaves, Persian clerks, Spanish judges, West African rulers, and Javanese performers within one us. Yet unresolved questions travelled with it from day one: what if people sharing commitments disagree, as beneath the shelter? What about neighbours with different commitments, as with Medina's Jewish tribes? What if individual conscience conflicts with communal commands, as for Ibn Hanbal?

These questions now pass to you, for they never belonged only to Muslims.

What binds you and your classmates into us? What should bind a class, team, country, or online community: common texts, rituals, interests, enemies, or something else?

Further, when the community's voice clashes with an individual's, who yields? Should Bilal fall silent for his master, or should Mecca's aristocrats dismantle their order for strangers? Should Ibn Hanbal submit to the caliph, or should the caliph acknowledge a boundary around conscience?

Give your answer and reasons. There is no standard answer. But your wording, your use of us, already mints a coin.
